% !PRINTED 401
\weberpage{401}

If $y$ is replaced by an associated number $\pm\varepsilon^k y$, then $\xi_1$
passes into $\xi_1 + k$, where $k$ is a positive or negative integer,
and one can therefore always choose $y$ among the associated numbers
in one way, so that

\weq{8}{0 \leqq \xi_1 < 1}

holds. The number $y$ so determined is then the \gesperrt{reduced
number}.

If we observe further that the equations (7) do not change
when $y$ is replaced by $-y$, it follows that:

3. \gesperrt{In a family of associated numbers there are always
two and only two reduced numbers, which differ only in
sign.}

From (7) there results, if $\xi_1$ satisfies the condition (8):

\weq{9}{0 \leqq \log\left|\frac{y}{y'}\right| < \log\frac{\varepsilon}{\varepsilon'},}

and from this

\weq{10}{1 \leqq \left|\frac{y}{y'}\right| < \frac{\varepsilon}{\varepsilon'},}

and conversely there follow from (10) again the equations (7) with
the condition (8).

Since we have moreover \gesperrt{assumed} that $y$ is to have a \gesperrt{positive}
norm, $y$, $y'$ have the same sign, and this sign
is positive if $x_2$ is positive. For then $y - y' = x_2\sqrt{D}$ is
positive, and since $|y| > |y'|$, both $y$ and $y'$ must be positive.

In the special case $x_2 = 0$, $y$ and $y'$ are positive if
we assume $x_1$ to be positive. The ambiguity still remaining
in 3. is thereby removed.

Since therefore $y$ and $y'$ are positive, it follows from (10):

\weq{11}{y'\varepsilon - y\varepsilon' > 0.}

From which, by (4) and (6):

\weq{12}{2aUx_1 - (T-bU)x_2 > 0, \qquad x_2 \geqq 0.}

And these conditions, which we call the \gesperrt{isolation conditions},
already imply that $y'$, and consequently
$N(y)$, is positive. For by (3) we have
\[
T > U\sqrt{D},
\]
and accordingly there follows from (12):
\[
U[2ax_1 + (b-\sqrt{D})x_2] > 0.
\]
Thus it is proved:
%CONT>>
